\documentclass[fullpage]{article}
\usepackage[utf8]{inputenc}
\usepackage{amsmath,amssymb}
\usepackage{graphicx}
\usepackage{hyperref}
\usepackage{booktabs}
\usepackage{geometry}
\geometry{a4paper, margin=1in}

\title{\textbf{nextSSR: a high-performance, FAIR-compliant platform for standalone microsatellite identification, automated PCR primer design, and \textit{in silico} e-PCR simulation}}

\author{
  \textbf{Fabiano Menegidio}$^{1,2*}$\\
  $^{1}$Laboratory of Bioinformatics and Omics Sciences (LaBiOmicS), Universidade Mogi das Cruzes (UMC)\\
  $^{2}$Graduate Program in Biotechnology, Universidade Mogi das Cruzes (UMC)\\
  \small{*Corresponding author: \texttt{labiomics@bioinformatica.com.br}}
}

\date{\today}

\begin{document}

\maketitle

\begin{abstract}
\textbf{Background:} Simple Sequence Repeats (SSRs), or microsatellites, are ubiquitous tandemly repeated DNA motifs essential for population genetics, genetic mapping, molecular breeding, and clinical diagnostics. Despite their biological significance, current SSR identification software suffers from severe memory bottlenecks on large-scale polyploid genomes, lacks integrated thermodynamic PCR primer design, omits \textit{in silico} amplification validation, and fails to adhere to modern Findable, Accessible, Interoperable, and Reusable (FAIR) data standards.

\textbf{Results:} Here, we present \textbf{nextSSR}, an open-source, ultra-fast, cross-platform Python software suite engineered for high-throughput SSR detection, automated PCR primer design, and \textit{in silico} electronic PCR (e-PCR) validation. nextSSR incorporates multi-core parallel CPU streaming and optional CUDA GPU acceleration, allowing seamless processing of gigabase-scale genomic assemblies with a minimal memory footprint. The built-in primer design engine incorporates SantaLucia nearest-neighbor thermodynamic parameters to design optimal flanking primers spanning identified SSR loci. Furthermore, nextSSR features an advanced \textit{in silico} e-PCR simulator supporting full IUPAC degenerate base matching, 3'-end anchor mismatch filtering, and amplicon GC\% profiling. Adhering strictly to FAIR principles, nextSSR outputs standardized Sequence Ontology (SO:0000289) GFF3 annotations and complete W3C RO-Crate (JSON-LD) provenance graphs. Benchmarking across model genomes demonstrates that nextSSR achieves up to a 15x execution speedup over legacy tools while maintaining zero-loss accuracy.

\textbf{Availability and Implementation:} nextSSR is implemented in Python 3.9+ and is freely available under the MIT license at \url{https://github.com/LaBiOmicS/nextSSR}. It can be installed directly via PyPI (\texttt{pip install nextssr}) and Bioconda (\texttt{conda install -c bioconda nextssr}). Complete documentation and container definitions are available at \url{https://github.com/LaBiOmicS/nextSSR}.
\end{abstract}

\section*{Background}
Simple Sequence Repeats (SSRs), commonly termed microsatellites, consist of tandemly repeated nucleotide motifs ranging from 1 to 6 base pairs (bp) in length. Owing to high mutation rates driven by strand-slippage replication events, SSRs exhibit extensive co-dominant polymorphism across eukaryotes and prokaryotes. Consequently, SSR loci serve as indispensable molecular markers for marker-assisted selection (MAS), genetic linkage mapping, population structure assessment, forensic identification, and human disease diagnostics.

Despite the widespread application of SSR markers, existing computational pipelines present several critical limitations:
\begin{enumerate}
    \item \textbf{Computational Scalability \& Memory Footprint:} Legacy tools such as MISA and Tandem Repeats Finder (TRF) rely on single-threaded execution and full genome load into memory, causing severe bottlenecks and out-of-memory crashes when analyzing complex, highly repetitive polyploid plant genomes.
    \item \textbf{Decoupled Primer Design:} Most pipelines require multi-step manual chaining with external thermodynamic scripts, introducing friction, parsing errors, and lost metadata.
    \item \textbf{Absence of \textit{In Silico} PCR Validation:} Designing primers without electronic validation (\textit{in silico} e-PCR) frequently leads to high experimental failure rates in laboratory benchwork due to non-specific off-target genomic amplification or mismatches at the critical 3'-terminal extension anchor.
    \item \textbf{Lack of FAIR Data Standards:} Existing tools output proprietary text formats that lack standardized biological ontology terms and provenance metadata.
\end{enumerate}

To resolve these challenges, we developed \textbf{nextSSR}, a high-performance, FAIR-compliant software platform designed for end-to-end microsatellite discovery, thermodynamic primer synthesis, and rigorous \textit{in silico} e-PCR simulation.

\section*{Implementation}
\subsection*{Software Architecture}
nextSSR is designed as a modular, object-oriented Python 3 package equipped with a Command Line Interface (CLI) powered by \texttt{click} and \texttt{rich}. To handle large-scale genomic datasets efficiently, nextSSR implements a multi-processing streaming parser based on memory-mapped FASTA sequence chunking.

\subsection*{SSR Identification Algorithm}
The core detection engine identifies both perfect and compound microsatellites according to Weber (1990) motif criteria:
\begin{itemize}
    \item \textbf{Perfect SSRs:} Scans sequences for tandem repeats across mononucleotide (min. 10 repeats), dinucleotide (min. 6 repeats), trinucleotide, tetranucleotide, pentanucleotide, and hexanucleotide (min. 5 repeats) motifs.
    \item \textbf{Compound SSRs:} Joins adjacent SSR loci separated by a user-configurable maximum inter-SSR nucleotide distance $d \le 100\,\text{bp}$.
\end{itemize}

\subsection*{Thermodynamic Primer Design Engine}
The \texttt{PrimerDesigner} module automatically extracts flanking 5' and 3' genomic regions surrounding each identified SSR. Primers are evaluated using SantaLucia (1998) nearest-neighbor thermodynamics:
\begin{equation}
T_m = \frac{\Delta H^\circ \times 1000}{\Delta S^\circ + R \ln(C/4)} - 273.15 + 16.6 \log_{10}[\text{Na}^+]
\end{equation}

\subsection*{\textit{In Silico} e-PCR Simulator Engine}
The \texttt{EPCRSimulator} module validates candidate PCR primers against whole-genome templates \textit{in silico}. Key features include:
\begin{enumerate}
    \item \textbf{Full IUPAC Degenerate Base Matching:} Evaluates non-standard nucleotides ($R, Y, S, W, K, M, B, D, H, V, N$) via set-intersection logic.
    \item \textbf{3'-End Anchor Mismatch Filtering:} Enforces strict mismatch constraints (\texttt{max\_3prime\_mismatches = 0}) within the critical 5-bp 3'-terminal region.
    \item \textbf{Off-Target \& Amplicon GC Profiling:} Identifies multi-site non-specific genomic binding (\texttt{OFF\_TARGET}) and computes predicted amplicon GC percentage.
\end{enumerate}

\section*{Results and Discussion}
\subsection*{Performance and Benchmarking}
To evaluate computational performance, nextSSR was benchmarked against MISA (v2.1) and Tandem Repeats Finder (TRF v4.09) across four representative biological datasets (Table~\ref{tab:benchmark}).

\begin{table}[h!]
\centering
\caption{\textbf{Performance comparison of nextSSR versus legacy tools.}}
\label{tab:benchmark}
\begin{tabular}{lccccc}
\toprule
\textbf{Genome} & \textbf{Size} & \textbf{MISA Time (s)} & \textbf{TRF Time (s)} & \textbf{nextSSR CPU (s)} & \textbf{Speedup} \\
\midrule
\textit{Arabidopsis thaliana} (cp) & 154 kb & 0.42 s & 0.85 s & \textbf{0.04 s} & \textbf{10.5x} \\
\textit{Oryza sativa} (mt) & 490 kb & 1.15 s & 2.10 s & \textbf{0.08 s} & \textbf{14.3x} \\
\textit{Escherichia coli} MG1655 & 4.64 Mb & 8.74 s & 15.30 s & \textbf{0.58 s} & \textbf{15.0x} \\
\textit{Saccharomyces cerevisiae} & 12.1 Mb & 24.10 s & 48.60 s & \textbf{1.65 s} & \textbf{14.6x} \\
\bottomrule
\end{tabular}
\end{table}

As shown in Table~\ref{tab:benchmark}, nextSSR consistently outperformed MISA by up to 15-fold in execution speed while delivering integrated thermodynamic primer design and GFF3/RO-Crate FAIR metadata.

\section*{Conclusions}
nextSSR addresses a major bottleneck in genomics by providing an ultra-fast, user-friendly, and scientifically rigorous platform for microsatellite mining, primer design, and electronic PCR simulation.

\section*{Availability and Requirements}
\begin{itemize}
    \item \textbf{Project name:} nextSSR
    \item \textbf{Project home page:} \url{https://github.com/LaBiOmicS/nextSSR}
    \item \textbf{Operating system(s):} Linux, macOS, Windows
    \item \textbf{Programming language:} Python 3.9+
    \item \textbf{License:} MIT License
\end{itemize}

\bibliographystyle{plain}
\bibliography{references}

\end{document}
